这里补一条衔接：前面的均方可导过程与布朗运动不能直接混同。
布朗增量 $W(t+h)-W(t)$ 的方差为 $h$，
所以差商的二阶矩为 $1/h$，在 $h\to0$ 时发散，布朗运动没有这样的均方导数。
因此 $\mathrm dW$ 只是增量与随机积分的记号，不能把它当作普通可微函数的微分来运算。

适应且平方可积的被积过程还满足 Itô 等距：
\[
\mathbb E\left|\int_0^T H_t\,\mathrm dW_t\right|^2
=\mathbb E\int_0^T|H_t|^2\,\mathrm dt.
\]
简单过程时，展开求和的平方，独立的未来布朗增量使交叉项期望为零，
每个对角项则给出 $\mathbb E[H_{t_i}^2]\Delta t_i$；再用均方极限推广。
这就是使用二阶矩构造随机积分的一条具体理由。
对 $\mathrm dX=\mu\,\mathrm dt+\sigma\,\mathrm dW$，
Itô 引理中的二阶项于是成为 $\tfrac12\sigma^2 f_{xx}\,\mathrm dt$，
下一节推导密度演化时出现的二阶空间导数正来自这里。
